% ==STYLE-BEGIN==
\documentclass[UTF8,11pt]{ctexart}
\usepackage{amsmath,amssymb,amsthm}
\usepackage[a4paper,margin=2.4cm,top=2.3cm,bottom=2.3cm]{geometry}
\usepackage{fancyhdr}
\usepackage{enumitem}
\usepackage{hyperref}
\hypersetup{hidelinks}

% ---- 版式 ----
\linespread{1.15}
\setlength{\parskip}{0.35em}
\setlength{\parindent}{1.6em}
\setlist{nosep,leftmargin=2em}

% ---- 页眉 ----
\pagestyle{fancy}
\fancyhf{}
\fancyhead[L]{\small\coursename}
\fancyhead[R]{\small\thepage}
\renewcommand{\headrulewidth}{0.4pt}

% ---- 定理环境：中文用正体，标题加粗，按小节共用一条流水号 ----
\newtheoremstyle{cnthm}{0.7em}{0.7em}{\normalfont}{}{\bfseries}{}{0.5em}{}
\theoremstyle{cnthm}
\newtheorem{thm}{定理}[section]
\newtheorem{lem}[thm]{引理}
\newtheorem{prop}[thm]{命题}
\newtheorem{cor}[thm]{推论}
\newtheorem{defn}[thm]{定义}
\newtheorem{rem}[thm]{注}
\newtheorem{exa}[thm]{例}
\newtheorem{conj}[thm]{猜想}
\renewcommand{\proofname}{\bfseries 证明}
\providecommand{\coursename}{}
% ==STYLE-END==
\renewcommand{\coursename}{数学分析\,·\,1015}
\title{数学分析 1015}
\date{}
\begin{document}
\maketitle

\section{Recall}
复合函数求极限
设$h(x) = f(g(x))$ , 令$u = g(x)$. 如果当$x \to x_0$时 , $g(x) \to u_0$且$\exists x_0$的去心邻域$\mathring{U}$满足$\forall x \in \mathring{U} , g(x) \neq U_0$,则
\[
\lim_{x \to x_0} f(g(x)) = \lim_{u \to u_0}f(u)
\]
\begin{quote}
大部分情况下(对于初等函数而言) , $\lim\limits_{u \to u_0} f(x) = f(u_0)$
\end{quote}

\begin{defn}
$\lim\limits_{x \to \infty} f(x) = A$是指$\forall \varepsilon >0 , \exists x , |x| > N$,有$|f(x)-A|< \varepsilon$
\end{defn}

命题
\[
\lim_{x \to x_0}f(x) = A \Leftrightarrow \lim_{x \to x_0^-}f(x) = \lim_{x \to x_0^+}f(x) = A
\]
我们已知趋向一个点时的情况,下面将其扩展到趋向无穷
\[
\lim_{x \to \infty}f(x) = A \Leftrightarrow \lim_{x \to +\infty}f(x) = \lim_{x \to -\infty}f(x) = A
\]

\section{无穷小量}
定义:若$\lim\limits_{x \to \infty} f(x) = 0$则称f(x)为$x \to x_0$的无穷小量

\begin{quote}
1. 无穷小量必须强调极限过程
\end{quote}
\begin{quote}
2. 无穷小量是函数而非常数0
\end{quote}

性质:设在同一极限过程中有$x \to x_0$
\[
\lim_{x \to x_0} f(x) = \lim_{x \to x_0} g(x) = 0
\]
则有
\noindent\textbf{1.} $f \pm g , f \times g$也是$x \to x_0$的无穷小量
\noindent\textbf{2.} 若$h$在$x_0$的某个去心邻域$\mathring{U}$中有界,则$h \times f$是$x \to x_0$时的无穷小量
\noindent\textbf{3.} 若$\lim\limits_{x \to x_0} \eta(x) = A$,则$n \times f$是$x \to x_0$时的无穷小量

\begin{quote}
3由局部有界性可以转换成2
\end{quote}

下对2进行证明
引理:
\[
\lim_{x \to x_0}f(x) = 0 \Leftrightarrow \lim_{x \to x_0}|f(x)| = 0
\]
在$\mathring{U}$处考虑
\[
\forall x \in \mathring{U} , 0 \leq |h(x)f(x)| \leq M|f(x)|
\]
由夹迫准则可知
\[
\lim_{x \to x_0} h(x)f(x) = 0
\]

\section{从无穷小量到泰勒级数}
注意到
\[
\lim_{x \to x_0}f(x) = A
\Leftrightarrow \lim_{x \to x_0}(f(x)-A) = 0
\Leftrightarrow f(x)-A \text{是无穷小量}
\]
不难发现
\[
f(x) = A + (f(x)-A)
\]
可以视作使用常数去近似$f(x)$

\begin{defn}
若$f,g$为$x \to x_0$时的无穷小量且$\lim\limits_{x \to x_0} \frac{f}{g} = 0$ , 则称f为g的高阶无穷小
\end{defn}

在上面的例子中,$f(x)-A$就是$A$的高阶无穷小
不妨更进一步,用多项式去近似$f(x)$,即
\[
f(x) = A + Bx + h(x)
\]
其中$h(x)$是x的高阶无穷小
由上不难发现
\[
\lim_{x \to 0}f(x) = A = f(0)
\]
考虑B,有
\[
\begin{aligned}
f(x) - A &= Bx +h(x) \\
\Rightarrow \frac{f(x)-A}{x} &= B + \frac{h(x)}{x} \\
\Rightarrow \lim_{x \to 0} \frac{f(x)-A}{x} &= B = \lim_{x \to 0} \frac{f(x)-f(0)}{x-0} = f(0)'
\end{aligned}
\]
再进一步
\[
f(x) = A +Bx +Cx^2 + \delta(x)
\]
其中$C$和$f(0)''$有关

\section{无穷大量}
定义: 
\[
\lim_{x \to x_0} f(x) = \infty \triangleq \forall N >0 , \exists \delta >0 , \forall x \ 0<|x-x_0|<\delta , |f(x)| > N
\]
则称f(x)为$x \to x_0$时的无穷大量
\begin{quote}
趋势可以是$x \to x_0 \quad x \to x_0^+ \quad x \to x_0^-$及$x \to \infty \quad x \to +\infty \quad x \to -\infty$六种中的任何一种
\end{quote}

\begin{quote}
无界函数和无穷大量

例如函数$f(x) = x \cos{x}$,取$x = 2n\pi$ , 有$f(x_n) = 2n \pi$ 可知
\[
\lim_{n \to +\infty}f(x_n) = +\infty
\]
显然,$f(x)$是无界函数

我们取$y_n = 2n\pi + \frac{\pi}{2}$,有$f(y_n) = 0$,根据海涅原理可知,$f(x)$不是$x \to x_0$时的无穷大量
\end{quote}

我们可以将$f(x)$的极限分为以下几种类型
\noindent\textbf{1.} 极限存在且为0(无穷小量)
\noindent\textbf{2.} 极限存在且不为0
\noindent\textbf{3.} 极限不存在,函数值趋向无穷大(无穷大量)
\noindent\textbf{4.} 极限不存在,函数值震荡

\section{无穷小(大)量之间的比较}
定义:
若$\lim\limits_{x \to x_0}f(x) = \lim\limits_{x \to x_0}g(x) = 0$且$\lim\limits_{x \to x_0} \frac{f(x)}{g(x)} = A \neq 0$,则称f与g同阶
如果$A = 1$ , 则称f与g是等价无穷小

我们有如下的常用结论
对于数列,
\[
\log{n}<n<n\log{n}<n^2<\cdots<n^k<\cdots<e^n<n!<n^n \quad n \to +\infty
\]
对于函数,
\[
\log{x}<x<x\log{x}<x^2<\cdots<x^k<\cdots<e^x
\]

对于等价无穷小,我们有
\[
\begin{aligned}
\ln{1+x} &\sim x\\
e^x - 1 &\sim x\\
a^x - 1 &\sim x\ln{a}\\
(1+x)^\alpha -1 &~ \alpha x
\end{aligned}
\]

\end{document}

\end{document}
